% ============================================================================
% 内存管理
% ============================================================================

\subsection{内存管理}

我们在裸机环境中实现了一个简单的内存分配器，支持分配和释放，还支持块合并。

内存布局：
- 每个内存块有 8 字节头部：
  - 4 字节：块大小
  - 4 字节：使用标记（0 未使用，1 已使用）
- 堆大小：4096 字节

内存分配器有三个函数：
\begin{itemize}
    \item \texttt{memory\_init}：初始化堆
    \item \texttt{memory\_alloc}：分配内存
    \item \texttt{memory\_free}：释放内存
\end{itemize}

\subsubsection{汇编实现}

\lstinputlisting[style=asm]{../QEMU/src/ASM/memory.s}

\subsubsection{C 语言封装}

C 语言中有对汇编函数的声明：

\lstinputlisting[style=cstyle]{../QEMU/src/include/memory.h}

\subsubsection{工作原理}
\begin{enumerate}
    \item \texttt{memory\_init}：初始化堆，第一个块大小设为 4088 字节（4096 - 8 字节头部）
    \item \texttt{memory\_alloc}：首次适配算法，找到足够大的空闲块
    \item \texttt{memory\_free}：释放内存并尝试与相邻空闲块合并
\end{enumerate}
